\section{Open Questions}
\label{sec:open-questions}

The following five questions are genuinely open after this replication: they are neither answered by the paper nor by our reproduction, and each has a concrete, low-cost next step.  They are grounded in a re-read of Cordoni (2023) and its GSM\textsuperscript{2} lineage (Cordoni \& Missiaggia \emph{et al.} 2021, \emph{Phys.\ Rev.\ E} 103:012412).

\subsection*{Q1.  Poisson vs.\ negative-binomial DSB induction at high LET}

\textbf{Question.}  Does the sub-Poissonian lethal-lesion result survive when the primary DSB induction distribution is over-dispersed (negative binomial, as expected at high LET), rather than treated as a deterministic initial condition $X_0$?

\textbf{Basis.}  The paper assumes $X_0$ is a fixed initial count and derives $\mathrm{Fano}(Y)<1$ from clustering + repair dynamics alone.  At high LET, primary DSB counts per cell are empirically and mechanistically over-dispersed --- spatially clustered on sub-nuclear scales (Friedland et al.\ 2017 track-structure work; Ballarini \& Ottolenghi 2005).  If the input $X_0\sim\mathrm{NegBin}(\mu,\phi)$ with $\phi\gg 1$, the output variance of $Y$ inherits an input-driven component that may swamp the paper's $-\delta(t)$ contraction.  The paper does not test this, and its ``sub-Poissonian'' headline could be misleading in the very regime (heavy ions, boron capture, proton Bragg peak) where LUCID users care most.

\textbf{Next steps.}  Rerun the existing SSA (\verb|code/gsm2_model.py|) with $X_0\sim\mathrm{NegBin}(\mu=100,\phi)$ for $\phi\in\{1,2,5,10\}$ instead of a fixed $X_0=100$.  Report $\mathrm{Fano}(Y)$ at $t=1.5$ as a function of $\phi$.  If $\mathrm{Fano}(Y)$ crosses~1, identify the $\phi_\mathrm{crit}$ at which the paper's headline flips.  Cost: $<1$~h CPU, \$0.  A strict extension of the existing single-parameter replication with a 5-line input distribution change.

\subsection*{Q2.  LNA in the high-clustering / small-$x_0$ regime}

\textbf{Question.}  How well does the linear-noise approximation hold in the high-clustering regime (large $\tilde b_K$ or small $r$), where the MME non-linearity dominates and higher-order van~Kampen corrections may become $\mathcal O(1)$?

\textbf{Basis.}  The paper's Sec.~4 numerics fix $(r,a,\tilde b_K)=(4,0.1,0.01)$, a regime where the clustering rate $\tilde b_K \bar x^2\le 100\cdot 0.01=1$ is small compared to $r\bar x\le 400$ over most of the trajectory.  The van~Kampen expansion is a leading-order asymptotic in $1/\sqrt K$; the paper does not quantify when the next-order correction becomes important.  LNA is known to fail (i) as $\bar x\to 0$ (paper's own caveat, our C4) and (ii) when the non-linearity in the drift is comparable to the diffusion coefficient.  Regime (ii) is unexplored.

\textbf{Next steps.}  Parameter sweep: fix $(a,r)=(0.1,4)$ and scan $\tilde b_K\in\{0.01,0.05,0.1,0.5,1.0\}$ at fixed $x_0=100$.  Compare SSA and LNA $\mathrm{Fano}(Y)$ at $t=1.5$ and the Kolmogorov--Smirnov distance between the SSA and LNA marginals of $Y$ at $t=0.7$.  Report the $\tilde b_K$ at which $\mathrm{KS}>0.05$ (a reasonable failure threshold).  Also run $x_0\in\{10,30,100,300\}$ at the paper's default rates to expose the small-$x_0$ quantized regime where LNA fails from below.  Cost: $<2$~h CPU, \$0.

\subsection*{Q3.  Identifiability of $(r,a,b)$ from real radiobiological data}

\textbf{Question.}  Do the paper's model rate constants $(r,a,b)$ have an identifiable mapping to any published radiobiological observable (foci-count histograms, cell-survival curves, complex-DSB fractions)?

\textbf{Basis.}  The paper picks $(r,a,\tilde b_K)=(4,0.1,0.01)$ ``for illustration'' with no reference to a fit against experimental data.  Three constants over one observable (survival at one dose) is underdetermined; the constants are also correlated in their effect on $\mathrm{Fano}(Y)$ (increasing $r$ reduces mean $Y$; increasing $b$ increases the clustering variance).  Without an identifiability analysis (Fisher information matrix, profile likelihoods, or a joint fit to a dose-response + foci-histogram dataset), the model is a plausibility argument rather than a predictive tool.  This is the gap between ``the paper's math is right'' (which we have verified) and ``the paper's model is useful for LUCID predictions'' (which is not established).

\textbf{Next steps.}  Two probes.  (a) Compute the Fisher information matrix of $\mathrm{Fano}(Y)$ with respect to $(\log r,\log a,\log \tilde b_K)$ via finite differences on the moment ODE (Eq.~16) and report its condition number --- if $\kappa>10^4$ the constants are effectively unidentifiable from $\mathrm{Fano}(Y)$ alone.  (b) Fit $(r,a,\tilde b_K)$ to a public $\gamma$-H2AX foci-count dataset (Rothkamm \& L\"obrich 2003 or Barnard et al.\ 2013 supplementary tables); report best-fit constants, 95\% profile-likelihood intervals, and residual QQ plot.  Cost: (a) $\sim 30$~min CPU, \$0; (b) 1--2~days for data acquisition + fit, \$0 if datasets are public.

\subsection*{Q4.  Two-phase repair kinetics and cell-to-cell heterogeneity}

\textbf{Question.}  How is the sub-Poissonian conclusion modified by realistic two-phase repair kinetics (fast NHEJ + slow HR/backup) or by cell-to-cell heterogeneity in the rate constants?

\textbf{Basis.}  The paper assumes a single time-constant sub-lethal repair rate $r$.  Real DSB repair is well-established as bi-exponential (Rothkamm \& L\"obrich 2003; Metzger \& Iliakis 1991): a fast NHEJ component ($t_{1/2}\sim 10$--$30$~min) and a slow component ($t_{1/2}\sim 4$--$8$~h).  Similarly $(r,a,b)$ are near-certainly heterogeneous across cells (cell-cycle phase, chromatin state, individual variation).  The paper's headline could either strengthen or weaken under these extensions: heterogeneity of $r$ across a population is a classical mechanism for over-dispersion (a mixture of Poissons is over-dispersed), so it might push population $\mathrm{Fano}(Y)$ above 1 even when each individual cell is sub-Poissonian.

\textbf{Next steps.}  Extension of the same SSA.  (a) Split $X$ into $X_\mathrm{fast}$ and $X_\mathrm{slow}$ with rates $(r_\mathrm{fast}=8, r_\mathrm{slow}=0.5)$ at a 70/30 initial split; rerun and compute $\mathrm{Fano}(Y)$.  (b) Cell-to-cell heterogeneity: draw $r_i\sim\mathrm{LogNormal}(\log 4,\sigma)$ for $\sigma\in\{0.1,0.3,0.5,1.0\}$ across 20\,000 independent cells; compute the population $\mathrm{Fano}(Y)$ at $t=1.5$.  Report the $\sigma$ at which the population Fano crosses~1.  Cost: 2--4~h CPU, \$0.  Interpretation: quantifies how much of the ``lethal lesions are sub-Poissonian'' story is a per-cell statement vs.\ a population statement.

\subsection*{Q5.  Tail probabilities, moment closure, and low-dose cell survival}

\textbf{Question.}  Does the moment closure at second order (LNA / linear FPE) systematically underestimate the tail probability $P(Y\gg\bar y)$, and if so, what does that imply for extrapolating from $\mathrm{Fano}(Y)$ to cell-survival (which is dominated by the tail $P(Y\ge 1)$) at low dose?

\textbf{Basis.}  The LNA gives a Gaussian marginal for $Y$ with mean $\bar y(t)$ and variance $c_{vv}(t)$.  At the paper's parameter set, $c_{vv}(1.5)\approx 7.65$ and $\bar y(1.5)\approx 11.26$; the Gaussian tail $P(Y>\bar y+3\sigma)\approx 0.00135$, whereas the true CTMC (a compound counting process with a clustering term) is known to have heavier tails than Gaussian for such reaction networks (Kurtz 1978; Anderson \& Kurtz 2015).  For radiobiology this matters: at low dose the interesting quantity is $P(Y=0)$ (cell survival), which the LNA computes from the Gaussian CDF at $0$.  If the true $P(Y=0)$ differs from the Gaussian prediction by even a factor of 2, cell-survival curves derived from the moment ODE alone will be miscalibrated at exactly the low-dose regime where LNT-vs-threshold debates live.

\textbf{Next steps.}  From the existing 20\,000 SSA paths, tabulate the empirical CDF of $Y$ at $t=1.5$ and compare against the Gaussian CDF with matched (mean, variance).  Report (a) KS distance, (b) empirical vs.\ Gaussian $P(Y=0)$ and $P(Y\ge 20)$, (c) Anderson--Darling $p$-value.  If $\mathrm{KS}>0.05$ or tail ratios $>2$, escalate to a saddle-point or system-size correction estimate of the tail.  If the discrepancy is significant, quantify how it propagates to a Poisson-of-lethal-lesions cell survival $S(D)=P(Y=0\mid \text{dose }D)$ and how much the LNA-derived survival curve differs from the SSA-derived one at low doses (0.1--1~Gy).  Cost: $\sim 1$~h CPU for the empirical tests, \$0; $\sim 4$~h for the survival extrapolation, \$0.
